%%%%%%%%%%%%%%%%%%%%%%%%%%%%%%%%%%%%%%%%%%%%%%%%%%%%%%%%%%%%%%%%%%%%%%%%%%%%%%%
\section{Status of Research}
\label{sec:status_of_research}

The classical van der Pauw method determines the sheet resistance of a thin film
from four-terminal measurements without requiring a Hall-bar geometry
\cite{VanDerPauw1958}. Its central advantage is the separation of current
injection and voltage sensing: ideally, the voltage contacts draw no current, so
contact voltage losses do not directly enter the measured sheet resistance. This
is particularly useful when contacts are large, poorly defined, or strongly
condition-dependent.

Contact effects are a major limitation when mobilities are extracted from
two-terminal gated thin-film measurements. At
metal--organic interfaces, injection depends on energetic alignment, interfacial
states, morphology, traps, and gate bias. Conventional transfer curves can thus
contain both channel and contact contributions \cite{Natali2012ChargeInjection}.
Bittle et al. showed that gate-dependent contact resistance can strongly
overestimate field-effect mobility when standard transistor equations are applied
directly to two-terminal data \cite{Bittle2016MobilityOverestimation}. A high apparent mobility is therefore not sufficient evidence for fast
transport in the semiconductor film.

Rolin et al. extended the van der Pauw concept to a gated organic accumulation
layer and showed that the gate-dependent sheet conductance can be used to extract
a long-range thin-film mobility with reduced contact influence
\cite{Rolin2017GVDP}. Related work has since adapted gated van der Pauw ideas to
organic mixed conductors and thin-film transistors in which contact resistance is
important \cite{Bonafe2021EgVDP,Lee2024GVDPIGZO,Lee2025VoltageDrivenGVDP}.
Together, these studies show that reliable mobility extraction requires a method
that distinguishes sheet transport from contact-related voltage losses.

For practical devices, this separation is meaningful only if the measurement
conditions are checked. The imposed current must be high enough to produce a
measurable voltage but low enough to remain close to linear response. Gate
leakage must remain small, the extracted mobility should be reproducible for
different currents, and the resistance of the full current-injection path should
be compared with the four-terminal sheet result. Slow trap relaxation and circuit
RC time constants can also make nominally static data depend on measurement
history. For this reason, the present work places the current-driven gated van der Pauw
measurement at the centre of the electrical characterization.
The four-terminal analysis is complemented by an
evaluation of the complete current-injection path, DLTS-type transient measurements, and a simple amplifier test in order to identify residual series and time-dependent effects.
